\chapter{Về Nhân Cách}
\markboth{Về Nhân Cách}{On Character}

\section{Học giỏi mà không tử tế thì chưa đủ.}
\begin{truyen}{Chưa đủ}{Not Enough}
\chuhoa{T}{hầy từng nói với lớp: ``Giỏi là đáng quý. Nhưng nếu giỏi mà làm người khác tổn thương, coi thường bạn yếu hơn, hoặc chỉ biết nghĩ cho mình, thì vẫn chưa đủ. Học giỏi mà không tử tế là một tài năng còn thiếu trái tim.''}
\textit{(The teacher once told the class: ``Being capable is valuable. But if you are capable and still hurt others, look down on weaker friends, or think only of yourself, it is still not enough. Skill without kindness is talent missing a heart.'')}
\end{truyen}

\begin{baihoc}
Tri thức đẹp nhất khi đi cùng sự tử tế.
\textit{(Knowledge is most beautiful when it goes with kindness.)}
\end{baihoc}

\begin{ghinhoanh}
\textbf{Short English to remember:}

Being smart is not enough; be kind too.

\end{ghinhoanh}

\begin{tuhocnhanh}
1. Vì sao tử tế làm cho tri thức có giá trị hơn? \textit{Why does kindness make knowledge more valuable?}

2. Em có đang vô tình làm ai tổn thương vì sự hơn thua không? \textit{Are you unintentionally hurting someone because of comparison?}
\end{tuhocnhanh}
\ngancach

\section{Trung thực với bài làm là trung thực với chính mình.}
\begin{truyen}{Không gian lận chính mình}{Honesty With Yourself}
\chuhoa{K}{hi có em muốn nhìn bài bạn, thầy nói: ``Em có thể qua mặt điểm số một lúc. Nhưng em không qua mặt được chính mình. Trung thực với bài làm là trung thực với trình độ thật của mình~--- để còn biết phải sửa ở đâu.''}
\textit{(When a student wanted to copy from a friend, the teacher said: ``You may fool the score for a moment, but you cannot fool yourself. Honesty in your work is honesty with your real level, so you know what to improve.'')}
\end{truyen}

\begin{baihoc}
Sự thật có thể làm em ngại một lúc; gian dối sẽ làm em yếu rất lâu.
\textit{(Truth may embarrass you briefly; dishonesty weakens you for a long time.)}
\end{baihoc}

\begin{ghinhoanh}
\textbf{Short English to remember:}

Honesty shows you where to grow.

\end{ghinhoanh}

\begin{tuhocnhanh}
1. Em có dám nhìn đúng năng lực của mình không? \textit{Do you dare to see your true level?}

2. Trung thực giúp việc học tiến bộ thế nào? \textit{How does honesty help learning improve?}
\end{tuhocnhanh}
\ngancach

\section{Người biết xin lỗi đang mạnh lên chứ không yếu đi.}
\begin{truyen}{Sức mạnh của lời xin lỗi}{The Strength to Apologize}
\chuhoa{C}{ó em làm bạn buồn nhưng cố im. Thầy nói: ``Nhiều người nghĩ xin lỗi là thua. Không. Người biết xin lỗi đang mạnh lên. Vì em ấy dám nhìn lỗi mình và dám sửa.''}
\textit{(A student hurt a friend but stayed silent. The teacher said: ``Many think apologizing means losing. No. The one who can apologize is becoming stronger, because they dare to face and correct their fault.'')}
\end{truyen}

\begin{baihoc}
Xin lỗi đúng lúc không làm em nhỏ đi; nó làm em trưởng thành hơn.
\textit{(An apology at the right time does not make you smaller; it makes you more mature.)}
\end{baihoc}

\begin{ghinhoanh}
\textbf{Short English to remember:}

An honest apology is a sign of strength.

\end{ghinhoanh}

\begin{tuhocnhanh}
1. Vì sao nhiều người khó nói lời xin lỗi? \textit{Why do many people find it hard to apologize?}

2. Em có đang nợ ai một lời xin lỗi không? \textit{Do you owe someone an apology?}
\end{tuhocnhanh}
\ngancach

\section{Tôn trọng người khác là dấu hiệu của người có giáo dục.}
\begin{truyen}{Điều thấy rất rõ}{A Visible Sign of Education}
\chuhoa{T}{hầy bảo: ``Người có giáo dục không chỉ thể hiện ở điểm cao hay lời nói hay. Dấu hiệu thấy rất rõ là cách em tôn trọng người khác~--- từ bạn học, thầy cô, đến cô lao công và chú bảo vệ.''}
\textit{(The teacher said: ``An educated person is not shown only by high grades or fine words. The clearest sign is how you respect others, from classmates and teachers to the janitor and security guard.'')}
\end{truyen}

\begin{baihoc}
Sự tôn trọng là thứ không tốn tiền nhưng nói rất rõ em là ai.
\textit{(Respect costs nothing, but it clearly shows who you are.)}
\end{baihoc}

\begin{ghinhoanh}
\textbf{Short English to remember:}

Respect is visible education.

\end{ghinhoanh}

\begin{tuhocnhanh}
1. Em có cư xử khác nhau với người “quan trọng” và “không quan trọng” không? \textit{Do you behave differently toward “important” and “unimportant” people?}

2. Điều đó nói gì về mình? \textit{What does that say about you?}
\end{tuhocnhanh}
\ngancach

\section{Biết ơn làm cho tri thức có chiều sâu.}
\begin{truyen}{Biết ơn thì học sâu hơn}{Gratitude Deepens Learning}
\chuhoa{T}{hầy nói: ``Một người chỉ biết nhận tri thức mà không biết ơn sẽ dễ thành kiêu. Còn người biết ơn thầy cô, cha mẹ, bạn bè, người giúp mình học được điều gì đó~--- sẽ có chiều sâu hơn trong cách sống.''}
\textit{(The teacher said: ``A person who only receives knowledge without gratitude easily becomes proud. But one who is thankful to teachers, parents, friends, and anyone who helps them learn becomes deeper in the way they live.'')}
\end{truyen}

\begin{baihoc}
Biết ơn giữ cho người học vừa sáng trí vừa ấm lòng.
\textit{(Gratitude keeps a learner both clear-minded and warm-hearted.)}
\end{baihoc}

\begin{ghinhoanh}
\textbf{Short English to remember:}

Gratitude gives depth to knowledge.

\end{ghinhoanh}

\begin{tuhocnhanh}
1. Ai đang góp phần vào việc học của em mỗi ngày? \textit{Who is helping your learning every day?}

2. Em đã từng nói lời cảm ơn với họ chưa? \textit{Have you thanked them?}
\end{tuhocnhanh}
\ngancach

\section{Giữ lời hứa nhỏ mới gánh được việc lớn.}
\begin{truyen}{Từ việc nhỏ}{From Small Promises}
\chuhoa{C}{ó em muốn làm lớp trưởng, muốn gánh việc lớn. Thầy hỏi: ``Em có giữ được lời hứa nhỏ như mang tài liệu, nộp bài, đến đúng hẹn chưa?'' Rồi thầy nói: ``Người gánh được việc lớn là người không xem nhẹ việc nhỏ.''}
\textit{(A student wanted to lead the class and handle bigger responsibilities. The teacher asked: ``Can you keep small promises like bringing documents, submitting work, and arriving on time?'' Then he said: ``The one who can carry big responsibilities is the one who does not treat small ones lightly.'')}
\end{truyen}

\begin{baihoc}
Năng lực lớn thường được chứng minh trước bằng sự chắc chắn ở việc nhỏ.
\textit{(Great ability is usually proven first through reliability in small things.)}
\end{baihoc}

\begin{ghinhoanh}
\textbf{Short English to remember:}

Big responsibility begins with small faithfulness.

\end{ghinhoanh}

\begin{tuhocnhanh}
1. Em đang xem nhẹ việc nhỏ nào? \textit{Which small responsibility are you taking lightly?}

2. Nếu sửa nó, em sẽ đáng tin hơn thế nào? \textit{How would fixing it make you more trustworthy?}
\end{tuhocnhanh}
\ngancach

\section{Đừng lấy thành công để coi nhẹ người khác.}
\begin{truyen}{Không dùng thành công để đứng cao hơn}{Success Without Arrogance}
\chuhoa{T}{hầy nói: ``Có chút thành công mà coi nhẹ người khác thì thành công ấy còn non. Thành công thật không cần em đứng lên vai ai để thấy mình cao. Nó làm em bình tĩnh, khiêm tốn và biết giúp người khác hơn.''}
\textit{(The teacher said: ``If you use a little success to look down on others, that success is still immature. Real success does not need you to stand on someone else's shoulders to feel tall. It makes you calmer, humbler, and more helpful.'')}
\end{truyen}

\begin{baihoc}
Càng có năng lực, em càng cần giữ sự khiêm tốn cho đẹp.
\textit{(The more capable you become, the more you need humility to remain beautiful.)}
\end{baihoc}

\begin{ghinhoanh}
\textbf{Short English to remember:}

Success should make you gentler, not higher.

\end{ghinhoanh}

\begin{tuhocnhanh}
1. Thành công nhỏ có đang làm em kiêu không? \textit{Is a small success making you proud?}

2. Em sẽ giữ sự khiêm tốn bằng cách nào? \textit{How will you protect humility?}
\end{tuhocnhanh}
\ngancach

\section{Nhân cách là thứ đi trước thành tích.}
\begin{truyen}{Đi trước}{What Comes First}
\chuhoa{T}{hầy bảo phụ huynh: ``Thành tích làm tôi mừng. Nhưng nhân cách mới làm tôi yên tâm. Vì thành tích có thể lên xuống; còn nhân cách mới quyết định một đứa trẻ sẽ sống thế nào khi không còn ai đứng cạnh.''}
\textit{(The teacher told the parents: ``Achievement makes me happy. But character gives me peace. Achievement can rise and fall; character decides how a child will live when no one is standing beside them.'')}
\end{truyen}

\begin{baihoc}
Nếu phải chọn nền trước, hãy chọn nhân cách trước rồi thành tích sẽ bền hơn.
\textit{(If you must choose a foundation first, choose character; achievement will then become more lasting.)}
\end{baihoc}

\begin{ghinhoanh}
\textbf{Short English to remember:}

Character comes before achievement.

\end{ghinhoanh}

\begin{tuhocnhanh}
1. Vì sao nhân cách bền hơn thành tích? \textit{Why is character more lasting than achievement?}

2. Em muốn người khác nhớ em bằng điều gì? \textit{How do you want people to remember you?}
\end{tuhocnhanh}
\ngancach

\section{Điều em làm khi không ai nhìn mới nói rõ em là ai.}
\begin{truyen}{Lúc không ai nhìn}{When No One Sees}
\chuhoa{T}{hầy nhặt một mẩu rác trong lớp rồi hỏi: ``Người vừa làm rơi có thể nghĩ không ai thấy. Nhưng điều em làm lúc không ai nhìn mới nói rõ em là ai. Nhân cách hiện ra mạnh nhất trong những lúc tưởng như rất kín đáo.''}
\textit{(The teacher picked up a piece of trash and asked: ``The one who dropped this may think no one saw it. But what you do when no one sees tells who you are most clearly. Character appears most strongly in moments that seem hidden.'')}
\end{truyen}

\begin{baihoc}
Nhân cách không chỉ là điều em nói về mình, mà là điều em làm trong im lặng.
\textit{(Character is not only what you say about yourself, but what you do in silence.)}
\end{baihoc}

\begin{ghinhoanh}
\textbf{Short English to remember:}

Who you are in private is who you are.

\end{ghinhoanh}

\begin{tuhocnhanh}
1. Nếu không ai nhìn, em có còn làm điều đúng không? \textit{If no one is watching, do you still do what is right?}

2. Điều đó cho em biết gì về mình? \textit{What does that tell you about yourself?}
\end{tuhocnhanh}
\ngancach

\section{Học để thành người trước khi thành tài.}
\begin{truyen}{Thành người trước}{Be Human First}
\chuhoa{T}{hầy kết quyển bằng câu quen thuộc: ``Học để thành tài là tốt. Nhưng học để thành người còn quan trọng hơn. Vì một người có tài mà thiếu nhân cách có thể gây hại rất lớn. Còn người tử tế sẽ biết dùng tài năng của mình đúng chỗ.''}
\textit{(The teacher closed the volume with a familiar line: ``Learning to become talented is good. But learning to become a good person matters even more. A talented person without character can do great harm. A kind person will know how to use talent rightly.'')}
\end{truyen}

\begin{baihoc}
Đích sâu nhất của giáo dục không chỉ là giỏi hơn, mà là tốt hơn.
\textit{(The deepest goal of education is not only to become more capable, but to become better.)}
\end{baihoc}

\begin{ghinhoanh}
\textbf{Short English to remember:}

Learn to be a good person before a successful one.

\end{ghinhoanh}

\begin{tuhocnhanh}
1. Em hiểu thế nào là “thành người”? \textit{What does it mean to “become a good person”?}

2. Em muốn dùng tri thức của mình để làm điều gì tốt? \textit{How do you want to use your knowledge for good?}
\end{tuhocnhanh}
\ngancach